% TOP_PROBLEM: {"id":30007589,"problem_number":"LOCAL-30007589","title":"Experimenter demand and elicitation","category_id":21,"category":{"id":21,"name":"economics","display_name":"Economics","description":"Open economic theory statements, published research agendas, and proposed scoped specifications; question provenance and historical priority are qualified per record.","slug":"economics","order_index":21,"created_at":"2026-10-08T00:00:00Z"},"status":"research_question","external_url":null,"legacy_ids":[],"published":false,"statement":"Identify or bound the neutral-demand treatment effect in an experiment by using randomized demand cues and the specified monotonicity and exclusion assumptions.","economics":{"original_id":78,"field":"Behavioral and experimental economics","kind":"identification","formalization_basis":"adapted_research_question","source_status":"research_agenda","priority_status":"earliest_found","priority_note":"The verified 2018 full manuscript directly proposes further study of demand mitigation. A 2017 working paper predates it; original priority of this longstanding methodological question is not established.","earliest_verified_reference":{"authors":["Jonathan de Quidt","Johannes Haushofer","Christopher Roth"],"title":"Measuring and Bounding Experimenter Demand","year":2018,"url":"https://jondequidt.com/pdfs/deQuidt_Haushofer_Roth_demand_final.pdf","doi":"10.1257/aer.20171330","locator":"Author manuscript dated April 4, 2018, §VI Conclusion, printed pp. 35–36 (PDF pp. 35–36), particularly last paragraph.","evidence_summary":"The paper derives demand-free bounds using randomized demand cues and explicitly proposes future work on mitigating demand and on how demand sensitivity changes with environmental features and incentives."},"current_reference":{"authors":["Jonathan de Quidt","Johannes Haushofer","Christopher Roth"],"title":"Measuring and Bounding Experimenter Demand","year":2018,"url":"https://jondequidt.com/pdfs/deQuidt_Haushofer_Roth_demand_final.pdf","doi":"10.1257/aer.20171330","locator":"Author manuscript dated April 4, 2018, §VI Conclusion, printed pp. 35–36 (PDF pp. 35–36), particularly last paragraph.","evidence_summary":"The paper derives demand-free bounds using randomized demand cues and explicitly proposes future work on mitigating demand and on how demand sensitivity changes with environmental features and incentives."},"formalization_note":"Adapted from the paper’s demand-bounding approach and explicit mitigation agenda. Not a claim that this established method remains undiscovered; the displayed bound uses explicitly stated additional assumptions."},"statusQualification":"Published question or research agenda verified at the cited locator. The mathematical specification is adapted for this catalogue and includes additional assumptions; source publication does not establish first-ever priority or certify this exact model remains unsolved. Adapted from the paper’s demand-bounding approach and explicit mitigation agenda. Not a claim that this established method remains undiscovered; the displayed bound uses explicitly stated additional assumptions.","statusReviewedAt":"2026-10-08"}
% FIRST_POSTED: 2026-10-08
% ENTRY_KIND: statement_only
% !TeX program = lualatex
\documentclass[11pt]{article}
\usepackage{fontspec}
\usepackage[a4paper,margin=25mm]{geometry}
\usepackage{amsmath,amssymb,mathtools}
\usepackage{xurl}
\usepackage[hidelinks]{hyperref}
\newcommand{\sref}[2]{\hyperlink{source:#1:#2}{[S#2]}}
\setlength{\emergencystretch}{3em}
\begin{document}
% problemId:30007589
\section*{Experimenter demand and elicitation}\label{30007589}
\subsection{Definitions and mathematical statement}
Let \(D\in\{0,1\}\) be an experimentally assigned substantive treatment and \(q\in[-1,1]\) a participant's belief about the action favored by the researcher, with \(q=0\) denoting a specified neutral-demand benchmark. Let \(Y_i(d,q)\) be the potential response of participant \(i\). Randomized demand cues \(Z\in\{-,0,+\}\) shift this belief while preserving the substantive task. The demand-free treatment estimand is
\[\tau^0=\mathbb E[Y_i(1,0)-Y_i(0,0)].\]
If negative and positive cues bracket neutral demand and \(Y_i(d,q)\) is monotone in \(q\), define \(\mu_d^-\) and \(\mu_d^+\) as mean outcomes under the negative and positive cues. Then the candidate bound is
\[\mu_1^- -\mu_0^+\leq\tau^0\leq\mu_1^+ -\mu_0^-.\]
Determine which cue and elicitation designs make these bracketing and exclusion restrictions credible across experimental environments, and which features reduce demand sensitivity without changing substantive preferences or information processing. Observed cue effects alone cannot verify that unobserved natural demand lies within the bracket. Measure cue comprehension, preregister effect-direction assumptions, and examine stronger cues, real incentives, anonymity, and survey versus laboratory contexts. Report partial identification when neutral behavior is not observed. The remaining agenda concerns validity and precision of separation in novel settings, rather than existence of a generic technique: such a bounding method is already supplied by the cited paper.

\par\textbf{Formulation and scope.} Adapted from the paper’s demand-bounding approach and explicit mitigation agenda. Not a claim that this established method remains undiscovered; the displayed bound uses explicitly stated additional assumptions.

\par\textbf{Open-question provenance.} An explicit question or research agenda was verified at the source locator. This does not by itself certify that every later version remains unresolved \sref{30007589}{1}.

\par\textbf{Historical priority.} The verified 2018 full manuscript directly proposes further study of demand mitigation. A 2017 working paper predates it; original priority of this longstanding methodological question is not established.

\subsection{Short English statement}
Identify or bound the neutral-demand treatment effect in an experiment by using randomized demand cues and the specified monotonicity and exclusion assumptions.

\subsection{Sources}
\begin{itemize}\raggedright
\item[\textnormal{[S1]}]\hypertarget{source:30007589:1}{} Jonathan de Quidt, Johannes Haushofer, Christopher Roth. 2018. Measuring and Bounding Experimenter Demand. Author manuscript dated April 4, 2018, §VI Conclusion, printed pp. 35–36 (PDF pp. 35–36), particularly last paragraph..
\url{https://jondequidt.com/pdfs/deQuidt_Haushofer_Roth_demand_final.pdf}.
DOI: 10.1257/aer.20171330.
{\small\textit{Earliest verified question/agenda reference; historical priority qualified below. The paper derives demand-free bounds using randomized demand cues and explicitly proposes future work on mitigating demand and on how demand sensitivity changes with environmental features and incentives.}}
\end{itemize}
\end{document}
